\section{Discussion}
\label{sec:discussion}

In this retrospective cohort of ED encounters with documented pain assessments, time to first pain reassessment was associated with a combination of clinical presentation, operational workflow, and insurance coverage.
Higher initial pain, ambulance arrival, and night shift were associated with faster reassessment, consistent with prioritization of clinically salient or higher-intensity presentations.
Lower triage acuity (higher ESI level) was associated with slower reassessment, which may reflect appropriate de-emphasis of repeat assessment among lower-acuity patients or differences in nursing and physician workflows by acuity zone.

\subsection{Insurance and monitoring intensity}

The most consistent disparity signal was insurance.
Medicaid patients experienced substantially slower reassessment than privately insured patients, with minimal attenuation after adjustment for demographics, severity, and workflow.
Within-acuity analyses suggested persistence in higher-acuity strata (ESI 1--2 and ESI 3), arguing against complete explanation by triage mix alone.
Medicare patients also had slower reassessment, though effect sizes were smaller than for Medicaid.

These findings align with broader literature on differential processes of care by payer status, but they highlight a distinct quality dimension---ongoing pain monitoring rather than analgesic prescribing alone.
Slower reassessment may reflect documentation practices, staffing patterns, patient advocacy differences, or unmeasured clinical complexity.
Because reassessment is a process measure tied to guideline-concordant care, persistent insurance gaps warrant quality-improvement attention and prospective validation.

\subsection{Race, ethnicity, and attenuation}

Race and ethnicity associations were smaller and more sensitive to adjustment.
Black patients had slower reassessment in minimally adjusted models that attenuated after severity and workflow adjustment.
Hispanic patients showed no independent association in the primary model.
Asian patients had slower reassessment in M4, but severe-pain and within-stratum estimates were imprecise.
We caution against over-interpreting race/ethnicity contrasts given overlapping social determinants, insurance mix, and documentation patterns not fully captured in administrative data.

\subsection{Workflow, severity, and temporal stability}

Workflow variables complemented clinical severity in explaining reassessment timing.
Ambulance and night-shift associations may reflect triage intensity, trauma protocols, or staffing models rather than patient preference alone.
The absence of a secular trend in 60-minute reassessment rates suggests that documented reassessment practice did not materially improve over the study's de-identified year window, despite national emphasis on pain management quality.

\subsection{Pathway analyses and interpretation}

Post-analgesic sensitivity analyses showed largely overlapping reassessment curves by race/ethnicity, implying that demographic patterns in the primary analysis are not driven principally by differences after treatment initiation.
Disposition pathway models (M4 plus admitted vs.\ home) were examined separately because disposition occurs downstream of initial pain documentation; these analyses inform mediation hypotheses but should not be interpreted as primary causal adjustments.
Inverse probability weighting for Medicaid vs.\ private insurance (propensity scores using pre-workflow covariates) can be reported in the supplement to assess robustness to measured confounding.

\subsection{Limitations}

This study has important limitations.
Data come from a single academic center via MIMIC-IV; generalizability to community EDs may be limited.
The cohort was restricted to acute pancreatitis and trauma diagnoses to ensure a population with expected pain documentation; findings may not extend to all ED presentations.
Pain scores and reassessment reflect documentation and workflow, not necessarily unobserved patient symptoms.
We could not fully account for analgesic timing, provider behavior, or social support.
Comorbidity measurement depended on linked ICD flags.
De-identified calendar years preclude direct mapping to policy interventions.
Other arrival mode was heterogeneous and retained only in supplemental analyses.

\subsection{Conclusions}

Pain reassessment timing in the ED reflects clinical presentation and operational context, with insurance-based differences that persist after multivariable adjustment and within triage-acuity strata.
Pain reassessment should be considered alongside analgesic outcomes in equity-oriented quality assessment.
Future work should validate these associations in multi-center cohorts, link reassessment timing to patient-reported outcomes, and evaluate interventions that standardize reassessment across payer groups.
